% ===== PRÉAMBULE CANONIQUE — à coller VERBATIM en tête de CHAQUE .tex =====
\documentclass[11pt,a4paper]{article}
\usepackage[utf8]{inputenc}\usepackage[T1]{fontenc}\usepackage{lmodern}
\usepackage[french]{babel}
\usepackage{amsmath,amssymb,mathtools}\usepackage{icomma}
\usepackage{siunitx}\sisetup{locale=FR}  % \qty{}{}, \si{}, \ang{}
\usepackage{xcolor}\usepackage{colortbl}
\usepackage{tikz}\usepackage{pgfplots}\pgfplotsset{compat=1.18}
\usepackage[siunitx,europeanresistors]{circuitikz} % circuits (élec.)
\usepackage[most]{tcolorbox}\usepackage{enumitem}\usepackage{array,booktabs}
\usepackage[a4paper,margin=2.2cm]{geometry}
\usetikzlibrary{arrows.meta,calc,angles,quotes,positioning,patterns,%
  backgrounds,shapes.geometric,decorations.pathmorphing,decorations.markings,intersections}
\definecolor{primary}{RGB}{222,49,99}\definecolor{primarydark}{RGB}{150,22,55}
\definecolor{defcol}{RGB}{0,114,178}\definecolor{methcol}{RGB}{52,92,148}
\definecolor{excol}{RGB}{0,135,90}\definecolor{attncol}{RGB}{200,40,55}
\tcbset{boxbase/.style={enhanced,breakable,boxrule=0.9pt,arc=2.5pt,
  left=3mm,right=3mm,top=2.3mm,bottom=2.3mm,fonttitle=\bfseries,coltitle=white}}
% --- boîtes TUTORIEL (noms EXACTS, ne pas renommer) ---
\newtcolorbox{cle}[1][]{boxbase,colback=defcol!6,colframe=defcol,
  title={\textsc{À retenir}\ifx\relax#1\relax\relax\else\textmd{ : \itshape #1}\fi}}
\newtcolorbox{methode}[1][]{boxbase,colback=methcol!7,colframe=methcol,sharp corners,
  title={\textsc{Méthode}\ifx\relax#1\relax\relax\else\textmd{ --- \itshape #1}\fi}}
\newtcolorbox{exemplebox}[1][]{boxbase,colback=excol!5,colframe=excol,
  title={\textsc{Exemple}\ifx\relax#1\relax\relax\else\textmd{ : \itshape #1}\fi}}
\newtcolorbox{piege}[1][]{boxbase,colback=attncol!6,colframe=attncol,
  title={\textsc{Piège}\ifx\relax#1\relax\relax\else\textmd{ : \itshape #1}\fi}}
\newtcolorbox{remarque}[1][]{boxbase,colback=black!4,colframe=black!45,
  title={\textsc{Remarque}\ifx\relax#1\relax\relax\else\textmd{ : \itshape #1}\fi}}
% --- boîtes EXERCICE / CORRIGÉ (noms EXACTS, auto-comptées) ---
\newtcolorbox[auto counter]{exo}[1][]{boxbase,colback=primary!4,colframe=primary,
  title={\textsc{Exercice}~\thetcbcounter\ifx\relax#1\relax\relax\else\textmd{ --- \itshape #1}\fi}}
\newtcolorbox[auto counter]{corrige}[1][]{boxbase,colback=excol!4,colframe=excol,
  title={\textsc{Corrigé}~\thetcbcounter\ifx\relax#1\relax\relax\else\textmd{ --- \itshape #1}\fi}}
\newcommand{\vv}[1]{\vec{#1}}\newcommand{\ds}{\displaystyle}
\newcommand{\R}{\mathbb{R}}\newcommand{\N}{\mathbb{N}}\newcommand{\Z}{\mathbb{Z}}
% =========================================================================
\begin{document}

\begin{exo}[deux rayons dans un prisme $\ang{30}$–$\ang{60}$–$\ang{90}$]
Un prisme en verre d'indice $n=\sqrt{2}$, entouré d'air, a pour section un triangle rectangle
$ABC$ (angle droit en $A$, $\ang{60}$ en $B$, $\ang{30}$ en $C$). Le rayon \textbf{1} entre
perpendiculairement à la face $AB$, le rayon \textbf{2} perpendiculairement à la face $AC$ ; tous
deux vont frapper l'hypoténuse $BC$.
\begin{center}
\begin{tikzpicture}[scale=1.5,>=Stealth]
  \coordinate (A) at (0,0); \coordinate (B) at (0,1.732); \coordinate (C) at (3,0);
  \fill[defcol!22] (A)--(B)--(C)--cycle;
  \draw[black!70,line width=1pt] (A)--(B)--(C)--cycle;
  \node[left,font=\scriptsize] at (A){$A$}; \node[above,font=\scriptsize] at (B){$B$};
  \node[right,font=\scriptsize] at (C){$C$};
  \node[above left,font=\scriptsize] at ($(B)+(0.05,-0.1)$){\ang{60}};
  \node[below,font=\scriptsize] at ($(C)+(-0.25,0.02)$){\ang{30}};
  \draw (0.15,0)--(0.15,0.15)--(0,0.15);
  \node[font=\scriptsize] at (0.55,0.55){$n=\sqrt2$};
  % rayon 1 : perpendiculaire a AB (horizontal), entre en (0,1.15), hypotenuse y=1.732-0.577x -> x=1.0
  \coordinate (M1) at (0,1.15); \coordinate (N1) at (1.0,1.15);
  \draw[->,line width=0.9pt] (-0.85,1.15)--(M1);
  \draw[line width=0.9pt,primary] (M1)--(N1);
  \node[left,font=\scriptsize] at (-0.85,1.15){\textbf{1}};
  \fill (N1) circle (0.7pt);
  % rayon 2 : perpendiculaire a AC (vertical), entre en (1.6,0), hypotenuse y=1.732-0.577*1.6=0.81
  \coordinate (M2) at (1.6,0); \coordinate (N2) at (1.6,0.808);
  \draw[->,line width=0.9pt] (1.6,-0.8)--(M2);
  \draw[line width=0.9pt,excol] (M2)--(N2);
  \node[below,font=\scriptsize] at (1.6,-0.8){\textbf{2}};
  \fill (N2) circle (0.7pt);
\end{tikzpicture}
\end{center}
\begin{enumerate}[label=\alph*)]
  \item Calcule l'angle limite verre/air du prisme.
  \item Montre (par la géométrie du triangle) que le rayon \textbf{1} frappe $BC$ sous $\ang{60}$ et
        le rayon \textbf{2} sous $\ang{30}$.
  \item Lequel des deux rayons \textbf{traverse} $BC$ (sort du prisme) ? Lequel subit une réflexion
        totale ?
\end{enumerate}
\end{exo}

\begin{exo}[le périscope à deux prismes]
Un périscope est réalisé avec \textbf{deux prismes} en verre ($n=1{,}5$) à section triangle rectangle
isocèle. La lumière entre horizontalement par le haut, subit une réflexion totale dans le premier
prisme, descend, puis une seconde réflexion totale dans le second prisme.
\begin{center}
\begin{tikzpicture}[scale=0.95,>=Stealth]
  % prisme haut
  \fill[defcol!22] (0.4,3.4)--(2.0,3.4)--(2.0,1.8)--cycle;
  \draw[black!70,line width=0.8pt] (0.4,3.4)--(2.0,3.4)--(2.0,1.8)--cycle;
  % prisme bas
  \fill[defcol!22] (2.0,1.0)--(2.0,-0.6)--(3.6,-0.6)--cycle;
  \draw[black!70,line width=0.8pt] (2.0,1.0)--(2.0,-0.6)--(3.6,-0.6)--cycle;
  \coordinate (P) at (1.35,2.75);
  \coordinate (Q) at (2.35,0.05);
  \draw[->,line width=1pt] (-1.0,2.75)--(P);
  \node[above,font=\scriptsize] at (-0.2,2.78){entrée};
  \draw[line width=1pt,primary] (P)--(Q);
  \draw[->,line width=1pt,primary] (Q)--(4.6,0.05);
  \node[below,font=\scriptsize] at (4.0,0.02){sortie};
  \fill (P) circle (1pt); \fill (Q) circle (1pt);
  \node[font=\scriptsize] at (1.25,3.15){P$_1$}; \node[font=\scriptsize] at (2.75,-0.25){P$_2$};
\end{tikzpicture}
\end{center}
\begin{enumerate}[label=\alph*)]
  \item Sur l'hypoténuse de chaque prisme, l'angle d'incidence vaut $\ang{45}$. Vérifie qu'il y a
        bien réflexion totale ($\hat{\imath}_{\text{lim}}$ verre/air $\approx\ang{41.8}$).
  \item Quelle est la direction du rayon de sortie par rapport au rayon d'entrée ?
  \item Pourquoi préfère-t-on des prismes à réflexion totale plutôt que de simples miroirs métalliques ?
\end{enumerate}
\end{exo}

\begin{exo}[l'angle d'acceptance d'une fibre optique]
Une fibre optique est formée d'un \textbf{c\oe ur} en verre ($n_1=1{,}5$) entouré d'une \textbf{gaine}
d'indice plus faible ($n_2=1{,}4$). La face d'entrée (plane) est en contact avec l'air ($n_0=1$). Un
rayon entre par cette face sous un angle $\theta$ et se réfracte dans le c\oe ur, puis frappe la paroi
c\oe ur/gaine.
\begin{center}
\begin{tikzpicture}[scale=1.0,>=Stealth]
  \fill[defcol!12] (0,-1.3) rectangle (6.2,1.3);
  \fill[defcol!30] (0,-0.8) rectangle (6.2,0.8);
  \draw[black!55,line width=1pt] (0,0.8)--(6.2,0.8);
  \draw[black!55,line width=1pt] (0,-0.8)--(6.2,-0.8);
  \draw[primary,line width=1.3pt] (0,-1.3)--(0,1.3);
  \node[font=\scriptsize] at (4.6,0){c\oe ur $n_1$};
  \node[font=\scriptsize] at (4.6,1.05){gaine $n_2$};
  \node[left,font=\scriptsize] at (0,-0.9){air $n_0$};
  \coordinate (M) at (0,-0.2);
  \draw[->,line width=1pt] (-1.5,-1.0)--(M);
  \draw[dash pattern=on 3pt off 2pt,black!55] (M)--(1.4,-0.2);
  \draw[line width=0.5pt] (0.65,-0.2) arc (0:33:0.65);
  \node[font=\scriptsize] at (0.98,0.02){$\theta$};
  \coordinate (P) at (1.7,0.8);
  \draw[line width=1.1pt,primary] (M)--(P);
  \draw[dash pattern=on 3pt off 2pt,black!55] (P)--(1.7,0.1);
  \fill (P) circle (1pt); \node[above,font=\scriptsize] at (P){$P$};
  \draw[line width=1.1pt,primary] (P)--(3.7,-0.8);
\end{tikzpicture}
\end{center}
\begin{enumerate}[label=\alph*)]
  \item Calcule l'angle limite $\hat{\imath}_{\text{lim}}$ à la paroi c\oe ur/gaine.
  \item Quel est, dans le c\oe ur, l'angle de réfraction \emph{maximal} à l'entrée pour que le rayon
        soit encore guidé (réflexion totale en $P$) ?
  \item Déduis-en l'angle d'entrée maximal $\theta_{\max}$ (« angle d'acceptance »).
\end{enumerate}
\end{exo}

\end{document}
